\documentclass[12pt]{article}
\usepackage{hyperref}

\usepackage{graphicx}

\usepackage{mathtools}
\DeclarePairedDelimiter\ceil{\lceil}{\rceil}
\DeclarePairedDelimiter\floor{\lfloor}{\rfloor}

\usepackage{cancel}

\usepackage{amsmath}
\usepackage{amsfonts}
\usepackage{amssymb}

\usepackage{xcolor}
\usepackage{listings}
\definecolor{darkgreen}{rgb}{0.0, 0.2, 0.13}
\lstset{language=C++,
                basicstyle=\ttfamily,
                keywordstyle=\color{blue}\ttfamily,
                stringstyle=\color{red}\ttfamily,
                commentstyle=\color{green}\ttfamily,
                morecomment=[l][\color{magenta}]{\#}
}

\usepackage{breqn}
\usepackage[]{algorithm2e}
\usepackage{physics}

\newcommand\fnurl[2]{%
  \href{#2}{#1}\footnote{\url{#2}}%
}

\newcommand{\Four}{\mathcal{F}}
\renewcommand{\(}{\left(}
\renewcommand{\)}{\right)}
%% \renewcommand{\{}{\left{}
%% \renewcommand{\}}{\right}}

\usepackage{natbib}
\usepackage{bibentry}
\bibliographystyle{plain}


\title{Multi-GPU peformance plan for rocFFT}
\author{}



\def\no#1{\nobreak{#1}} % stop equation breaking in dmath

\begin{document}
\maketitle

\section{Introduction}

These are notes on how to achieve performance for multi-device FFTs with
rocFFT.

\section{Strategy}

The majority of user feedback for rocFFT is from users in a cluster
environment.  As such, and particularly with introduction of cuFFTMp,
NVIDIA's distributed FFT library, there is interest in developing a
distributed version of rocFFT.

Many cluster applications which use FFTs developed their own
distributed version.  For example, GESTS and gromacs.  In order for a
distributed rocFFT to be succesful, it needs to be faster / more
reliable / simpler to use / more feature complete than the application
implementation.  That said, we have had interest from several
applications for an AMD version.

cuFFTMp uses NVSHMEM for communication.  ROCSHMEM is an AMD research
project which doesn't have release-level support.  As such, MPI is
probably the best solution for a distributed rocFFT.

Distributed FFTs need to have data transposed so that 1D FFTs can be
computed with all the relevant data on a single device.  This
transpose is very expensive, and distributed FFTs are severely
bandwidth-limited.  Optimizing this stage is the key to producing a
good distributed FFT.  In section~\ref{section_data_decomp}, we
introduce the standard data decompositions and the associated
transpose costs.

In addition, we can look at using MPI sub-communicators, FFTW's
various algorithms for transpose, and FFTW++'s recursive transpose.
In addition, it may be possible to do some passes on distributed data
in order to get as much work done as possible when the data is
contiguous, as well as allowing us to separate data dependencies so as
to allow for computation/communication overlap.  See
section~\ref{subpasses}.

Compressing data may be a useful strategy.  zfp is available on the
GPU.  Lossy compression may also be appropriate if we can effectively
control the loss of precision.  In addition, the fast multipole method
may be worth considering, as this has generally better communication
properties than a transpose.  However, this requires the user to
specify a precision tolerance.

Since one of the main uses of the FFT is for convolutions (eg GESTS,
gromacs), we can design an distributed FFT that is faster for this
operation.  At the most basic level, FFT-based convolutions need to be
zero-padded.  For the normal $1/2$ padding case, a 3D convolution
requires a buffer that is $87.5\%$ zeros; the first transpose can skip
the $75\%$ of the buffer that is just zero.  Furthermore, implicitly
dealiased convolutions can further reduce computation cost and allow
for better communication patterns.



\section{Data Decompositions}
\label{section_data_decomp}


\subsection{2D Decompositions}

There are basically two decompositions for 2D data of size $N^2$ over
$P$ procs; either block or pencil.  Block decompositions have
dimensions
\begin{dmath}
  \label{brick2d}
  \left(\frac{N}{P^{1/2}}, \frac{N}{P^{1/2}} \right)
\end{dmath}
and pencil has dimensions
\begin{dmath}
  \label{pencil2d}
  \left(N, \frac{N}{P} \right).
\end{dmath}

\subsubsection{Brick-pencil}
Each of the $P$ procs has to send $\sqrt{P}-1$ messages with dimension
\begin{dmath}
  \left(\frac{N}{P^{1/2}}, \frac{N}{P} \right)
\end{dmath}
for a message size of $N^2/P^{3/2}$.  The total communication size is
\begin{dmath}
  P \(\sqrt{P}-1\) N^2/P^{3/2} \approx N^2
\end{dmath}

\subsubsection{Pencil-pencil}
Each of the $P$ procs has to send $P-1$ messages with dimension
\begin{dmath}
  \left(\frac{N}{P}, \frac{N}{P} \right)
\end{dmath}
for a message size of $N^2/P^2$.  The total communication size is
\begin{dmath}
  P \(P-1\) N^2/P^2\approx N^2.
\end{dmath}


\subsubsection{2D transform summary}

\begin{centering}
  \begin{table}[htbp]
    \begin{tabular}{|l|l|l|}
      \hline
      Type          & Message size  & Messages per proc \\
      \hline
      Brick-pencil  & $N^3/P^{3/2}$ & $\sqrt{P}$          \\
      Pencil-pencil & $N^3/P^2$ & $P$        \\
      \hline
    \end{tabular}
  \end{table}
\end{centering}




\subsection{3D Decompositions}

\url{https://www.osti.gov/servlets/purl/1483229}

Given data of size $N^3$ with $P$ mpi procs, the brick layout has
each process having
\begin{dmath}
  \label{brick3d}
  \left(\frac{N}{P^{1/3}}, \frac{N}{P^{1/3}}, \frac{N}{P^{1/3}} \right)
\end{dmath}
and this must be transformed into either a pencil decomposition
\begin{dmath}
  \label{pencil3d}
  \left(N, \frac{N}{\sqrt{P}}, \frac{N}{\sqrt{P}} \right)
\end{dmath}
or a slab decomposition,
\begin{dmath}
  \label{slab3d}
  \left(N, N, \frac{N}{P} \right)
\end{dmath}

\subsubsection{Brick-Pencil}
Going from distribution~\eqref{brick3d} to~\eqref{pencil3d}, each of
the $P$ procs need to receive $P^{1/3}-1$ messages of dimensions 
\begin{dmath}
  \left(\frac{N}{P^{1/3}}, \frac{N}{\sqrt{P}}, \frac{N}{\sqrt{P}} \right)
\end{dmath}
for message size of $\frac{N^3}{P^{4/3}}$.  There are therefore
$P(P^{1/3}-1)$ messages, for a total communication size of
\begin{dmath}
  \frac{N^3}{P^{4/3}} P(P^{1/3}-1) \approx N^3.
\end{dmath}

\subsubsection{Brick-Slab}
Going from distribution~\eqref{brick3d} to~\eqref{slab3d}, each of the
$P$ procs needs to receive $P^{2/3}-1$ messages of dimension
\begin{dmath}
  \left(\frac{N}{P}, \frac{N}{P^{1/3}}, \frac{N}{P^{1/3}} \right)
\end{dmath}
which has message size $\frac{N^2}{P^{5/3}}$.  The total communication
  size is
\begin{dmath}
  P\(P^{2/3}-1\) \frac{N^2}{P^{5/3}} \approx N^3.
\end{dmath}

\subsubsection{Pencil-Pencil}
A pencil-to-pencil transpose requires all $P$ procs to send
$\sqrt{P}-1$ messages of size
\begin{dmath}
  \left(\frac{N}{\sqrt{P}}, \frac{N}{\sqrt{P}}, \frac{N}{\sqrt{P}} \right)
\end{dmath}
which has message size $\frac{N^3}{P^{3/2}}$.  The total communication
size is
\begin{dmath}
  P \(\sqrt{P}-1\) \frac{N^3}{P^{3/2}} \approx N^3.
\end{dmath}

\subsubsection{Pencil-Slab}
All $P$ procs send $P-1$ messages of size
\begin{dmath}
  \left(\frac{N}{P}, \frac{N}{\sqrt{P}}, \frac{N}{\sqrt{P}} \right)
\end{dmath}
for a message size of $N^3/P^2$.  The total communication size is
\begin{dmath}
  P \(P-1\) \frac{N^3}{P^2} \approx N^3.
\end{dmath}

\subsubsection{Slab-Slab}
All $P$ procs send $P-1$ messages of size
\begin{dmath}
  \left(\frac{N}{P}, \frac{N}{P}, N \right)
\end{dmath}
for a message size of $N^3/P^2$. The total communication size is
\begin{dmath}
  P \(P-1\) \frac{N^3}{P^2} \approx N^3.
\end{dmath}

\subsubsection{3D transform summary}

\begin{centering}
  \begin{table}[htbp]
    \begin{tabular}{|l|l|l|}
      \hline
      Type          & Message size  & Messages per proc \\
      \hline
      Brick-pencil  & $N^3/P^{5/3}$ & $P^{1/3}$          \\
      Brick-slab    & $N^3/P^{5/3}$ & $P^{1/3}$          \\
      Pencil-pencil & $N^3/P^{3/2}$ & $P^{1/2}$        \\
      Pencil-slab   & $N^3/P^2$    & $P$               \\
      Slab-slab     & $N^3/P^2$    & $P$               \\
      \hline
    \end{tabular}
  \end{table}
\end{centering}

\section{Performance Strategy}

\subsection{Communication library}

The choice of communication library is quite important here.  Major
MPI libraries have been GPU-aware for some time now; these need to be
benchmarked.  In particular, \texttt{MPI_Alltoall},
\texttt{MPI_Alltoallv}, and \texttt{MPI_Alltoallw} (and there
non-blocking versions).  In particular, it seems that
\texttt{MPI_Alltoallw} lacks high-performance implementations.  In
addition, the scheduling algorithm
in~\cite{schreuder1980constructing}

Different types of GPU-aware MPI:
\begin{itemize}
\item stream triggered
\item kernel triggered
\item kernel initiate
\end{itemize}
FIXME: explain and links to documentation.

rocSHMEM is currently experimental; it needs users in order to be in
an official release - being in an official release would really help
increase the user base!  We also need to verify rocSHMEM support and
how MPI interoperatbility works.

RCCL is fully supported and offers some useful features such as
\texttt{ncclAllToAll}, which can be used in at least single-node
communications.

Unknowns:
\begin{itemize}
\item nvSHMEM performance relative to MPI
\item rocSHMEM release plan
\item rocSHMEM/MPI interoperability
\end{itemize}

\subsection{1D/2D decomposition}

Also known as slab and pencil.

We need to
\begin{itemize}
\item support both
\item Compute the communication cost for each of these (depending on
  latency and bandwidth).
\item be able to switch from one to another internally when
  tuning/estimation indicates a performance gain.
\end{itemize}


\subsection{Compression / intermediary precision}

We know, at the FFT-library level, a lot about our data.  It's an
array of complex (or sometimes real) values, and we are willing to
sacrifice some small amount of numerical error to gain performance -
the FFT, while highly accurate, still has roundoff error.  We have
opportunities to reduce transform time by compressing data before
communication.  It may be possible to do this when the data is already
in LDS, which would reduce the cost of compression.  Application-level
kernels may also be able to consume compressed data, which would be an
interesting addition to the API.

In addition, we may be willing to use a lower precision when
transferring over the network - FFTs could be done in double
precision, but transferred in single, for example.

All of this is fairly application-dependant.  Precision requirements,
numerical ranges, and numerical accuracy requirements will be
determined by the application developer and not the library developer.
As such, any lossy method needs to be user-controlled, and therefore
included in the API.

\subsection{De-distributing the transform}

FFTs are bandwidth-limited; distributed transforms even more so.  Some
applications may distribute the data to optimize for other parts of
the program; we are not necessarily given an FFT-optimized data
layout.  This is particular when the size of the data is small
relative to the size and number of devices.

If the transform is really small, it may make more sense to gather the
data onto one device and distribute the data via a scatter operation.
In order to further remove a communication stage, we ould also
all-gather, compute redundantly, and then just pass the relevant data
to the local device.  This can also be done in an intermediate
fashion, for example computing slabs from a pencil decomposition.

We can compute the cost of these based on latency and bandwidth and
compare with the slab/pencil decomposition costs.

\subsection{Sub-passes}
\label{subpasses}

In the case where one of the distributed dimensions has the same
number of values (in the dimension), we can perform partial passes in
that direction.  This allows us to perform some more computation while
we have the data locally, which may allow us to eliminate some MPI
transposes entirely.  Also of interest is that this will allow us to
de-couple the data dependance; for example, if we do a radix-2 pass,
then we now have two different data sets on which we can compute the
transform independently.  Then, we can send some data, start
computing, and send the second data set while we work on the first.
This will also allow for a non-blocking interface.  This is scoped in
single-device ticket
\href{https://ontrack-internal.amd.com/browse/SWDEV-392917}{SWDEV-392917}.



\subsection{Latency/bandwidth tradeoff}

Under a strong-scaling regime, transforms are eventually
latency-bound.  At this point, bandwidth and flops are kind of free,
and we can spend some of these to try and reduce the cost of latency.
This is the idea behind the recursive MPI transpose
in~\cite{bowmanroberts2015adaptive}.  This needs validation (and
publication).


\section{Spectral Operations}
\label{spectralops}

FFTs are generally used in a few operations such as convolution,
Poisson solve, filtering, etc.  The combination of spectral ops with
multi-device transforms allows for optimisation which would be quite
cumbersome for the user.  For example, a Poisson solve can be done as
a round-trip in the last FFT kernel, thus avoiding two global-memory
round-trips.  Convolution/correlation are even more interesting
possibilities, as the zero-padding increases the effective problem
size and therefore communication costs.  Using implicit dealiasing
(\cite{dealias2} allows for the overlapping of communication and
computation.

The spectral operations API should be developed in conjunction
with the multi-device API.

\section{Other libraries}

\subsection{cuFFT and other back-ends}

\texttt{cuFFTMp} is already available, and hipFFT should somehow wrap
this, while also allowing users to call rocFFTMp.  We plan on adding
new functionality to rocFFT/Mp; features not implemented in cuFFT will
return a not-implemented error code.


\subsection{Other ROCm libs}

Also of note is that cuSolverMP is a thing; rocSolver developers may
wish to participate in choices in rocFFTMp design so that the ROCm libraries
have some sort of coherence.


\section{Implimentation plan}

First, get it working.  Just move everything to pencil/slab, send/receive
to whatever crazy data distribution that they specify.

Optimizations:
\begin{itemize}
\item gather/scatter
\item sub-passes in discontiguous direction
\item Do not use \texttt{MPI_Alltoallw} until there is a reasonable
  implementation.  It might be worth just using \texttt{MPI_Alltoallv}
  and handling the data formatting on our own.
\item Can we do \texttt{MPI_Alltoallw} via \texttt{MPI_Alltoallv} +
  extra isend/irecv?
\item \dots
\end{itemize}



\bibliography{../refs}

\end{document}
